% questions.tex —— 1977 年·江苏农垦系统初试卷 · 题目内容片段（无答案）
% 由 tools/generate_projects.py 依来源总账自 1977/1977_数学真题_江苏农垦系统初试_试题.md 生成，勿手改（改源 MD 后重跑）


\gkpart{江苏农垦系统初试卷}

\begin{gktip}
\textbf{来源说明}：本题面依据固定源 PDF《1977江苏农垦系统初试.pdf》第 1—2 页逐字转录。原卷为 1977 年江苏农垦系统高校招生初试（预选）数学卷，页面顶部唯一标题为汇编分节标题「十四、江苏省农垦系统」（「十四」为汇编节序号，非试卷自身编号），本文件 H1 按其逐字转写；原卷题号为中文数字「一、」至「八、」，本库按规范转写为「\textbf{N．}」。全卷共 8 题，无分值标注、无附图、无附加题或参考题；原卷未附答案，配套答案与解析为编者独立推导，\textbf{非官方标准答案}。
\end{gktip}

\par\medskip

\begin{question}
\gkti{1}计算：

（1） $\text{〔}1\frac{1}{24}-\left(\dfrac38+\dfrac16-\dfrac34\right)\times24\text{〕}\div(-5)^0\text{\text{；}}$

（2） $\lg\dfrac52+\lg\dfrac25-\lg10^3\text{.}$

\end{question}
\begin{question}
\gkti{2}解方程

\nopagebreak[4]
\[
\frac1{x+2}+\frac{4x}{x^2-4}-\frac2{x-2}=1.
\]

\end{question}
\begin{question}
\gkti{3}用长 $60$ 米的篱笆围成一矩形养鸡场地，一面利用原有旧墙，问矩形的长和宽各多少米才能使所围成的面积最大？

\end{question}
\begin{question}
\gkti{4}把极坐标方程 $\rho=2\cos\theta$ 化成直角坐标方程，并说出它是什么曲线？

\end{question}
\begin{question}
\gkti{5}在锐角 $\triangle ABC$ 中，已知 $\sin A=\dfrac35$，$\cos B=\dfrac5{13}$，求 $\cos C$ 的值。

\end{question}
\begin{question}
\gkti{6}三个数成等差数列，且它们的和是 $9$，如果把这个数列依次加上 $1$、$1$、$3$，那么所得的数又成等比数列。求这三个数。

\end{question}
\begin{question}
\gkti{7}求以椭圆 $\dfrac{x^2}{49}+\dfrac{y^2}{25}=1$ 的焦点和长轴的端点作为顶点和焦点的双曲线方程。

\end{question}
\begin{question}
\gkti{8}等腰三角形 $ABC$ 的顶角 $A$ 为 $100^\circ$，底角 $B$ 的平分线交腰 $AC$ 于 $D$，求证：$BC=AD+BD$。

\end{question}